\documentclass{article}
\usepackage{amsmath,amssymb,xurl}
\usepackage[hidelinks]{hyperref}
% Shared commands for notes in docs/paper-gaps/.

\DeclareUnicodeCharacter{2080}{\ifmmode{}_{0}\else\textsubscript{0}\fi}
\DeclareUnicodeCharacter{2081}{\ifmmode{}_{1}\else\textsubscript{1}\fi}
\DeclareUnicodeCharacter{2082}{\ifmmode{}_{2}\else\textsubscript{2}\fi}
\DeclareUnicodeCharacter{2099}{\ifmmode{}_{n}\else\textsubscript{n}\fi}

\newcommand{\Lean}{\textsc{Lean}}
\ExplSyntaxOn
\NewDocumentCommand{\leanid}{m}{
  \ifmmode
    \textsf{\detokenize{#1}}
  \else
    \group_begin:
    \urlstyle{sf}
    \tl_set:Nn \l_tmpa_tl {#1}
    \tl_replace_all:Nnn \l_tmpa_tl {\_} {_}
    \exp_args:NV \path \l_tmpa_tl
    \group_end:
  \fi
}
\ExplSyntaxOff
\providecommand{\tr}{\operatorname{tr}}
\newcommand{\tnleanIssueBase}{https://github.com/LionSR/TNLean/issues/}
\newcommand{\tnleanPrBase}{https://github.com/LionSR/TNLean/pull/}
\newcommand{\tnleanDocBase}{https://sirui-lu.com/TNLean/docs/}
\newcommand{\ghissue}[1]{\href{\tnleanIssueBase#1}{GitHub issue \##1}}
\newcommand{\ghpr}[1]{\href{\tnleanPrBase#1}{PR \##1}}
\newcommand{\doclink}[1]{\href{\tnleanDocBase#1.html}{\path{#1}}}

% Dirac notation and the identity operator, as in blueprint/src/macros/common.tex.
% \providecommand keeps note-local definitions authoritative.
\usepackage{dsfont}
\providecommand{\ket}[1]{|#1\rangle}
\providecommand{\bra}[1]{\langle #1|}
\providecommand{\braket}[2]{\langle #1 | #2 \rangle}
\providecommand{\ketbra}[2]{|#1\rangle\!\langle #2|}
\providecommand{\Id}{\mathds{1}}
\providecommand{\one}{\mathds{1}}
\providecommand{\C}{\mathbb{C}}
\providecommand{\MN}[1]{M_{#1}(\C)}
\providecommand{\MD}{M_D(\C)}

% Verdict marker: \gapnote{<kind>}{<status>}. Kind is the policy
% classification (clarification, local-correction, scope-restriction,
% unfaithful, false-source, open-gap); status is open, wip (elimination
% underway), resolved, or historical. Machine-read by the site index;
% prints a verdict line.
\newcommand{\gapnote}[2]{\par\noindent\textbf{Verdict:} #1 (#2).\par}

\title{Vacuum Density for Central Complementary Insertions}
\author{}
\date{}
\begin{document}
\maketitle
\gapnote{scope-restriction}{open}

Schuch, Cirac, and P\'erez-Garc\'ia,
\href{https://arxiv.org/abs/1001.3807v3}{arXiv:1001.3807v3},
Definition 6.1 and Lemma 6.14, source
\path{Papers/1001.3807/paper_v3.tex}, lines 1692--1700 and 2199--2214,
provide the regular isometry and local flux-string claim.

For a finite native graph, take regular $G$-isometric original tensors
and a central group label $u_e$ on every oriented edge. Centrality makes
left multiplication by each $u_e$ commute with the common translation
appearing in the original local Gram kernel. The local twisted Gram
kernel is consequently identical to the original one. The actual
open-region expansion gives
\leanid{TNLean.PEPS.regularTwistedOpenRegionMatrix_gram_eq_of_central}
for every finite region, without connectivity or supplied Gram identities.

If $u_e=1$ on all edges incident to an observation region $R$, the
observed open-region contraction is unchanged. The complementary
open-region Gram matrix is unchanged by centrality. Substitution into
the actual physical cut factorization proves
\leanid{TNLean.PEPS.regularPhysicalCutMatrix_reduced_eq_of_central}:
\begin{align}
    M_uM_u^\dagger=M_0M_0^\dagger.
    \label{eq:central_flux_reduced}
\end{align}
The normalized identity is
\leanid{TNLean.PEPS.regularPhysicalCutMatrix_normalized_reduced_eq_of_central}.
These identities supply neither a nonzero-cut premise nor a desired
factorization or density equality as an assumption. The unnormalized
identity is meaningful even for zero cuts; nonvanishing of actual torus
flux cuts is available separately from regular injectivity.

This proves the complementary-contraction step for central insertions,
including all insertions in an abelian group. Identity on incident bonds
is an explicit structural hypothesis, not inferred from endpoint
exclusion. The physical detour theorem computes a sufficient forbidden
corner-collar geometry separately. Combining it with homotopy completeness
and the native-to-graph cut conversion is a further composition step.
Noncentral open flux pairs remain under \ghissue{8703}.

The annular Wilson-loop obstruction is consistent with
(\ref{eq:central_flux_reduced}): a path connecting an enclosed endpoint
to an exterior endpoint cannot make every incident annular bond insertion
identity while preserving its flux. Endpoint exclusion alone is therefore
insufficient. The full-lattice Hamiltonian scope of \ghissue{8947} is unchanged.
\end{document}
